\section{Data}
\label{sec:data}

\subsection{Sources}

\begin{tabularx}{\textwidth}{@{}>{\raggedright\arraybackslash}p{5cm}X@{}}
\toprule
Source & What it provides \\
\midrule
CRSP daily stock file (CIZ format, release \texttt{ciz202512}) & daily returns including delisting
returns, prices, open/high/low/close, closing bid and ask, volume, shares outstanding, market
capitalisation, index membership spells, index returns \\
CRSP/Compustat Merged (release \texttt{cfz202607}) & quarterly income statement, balance sheet and
cash flow; report dates (RDQ) and filing dates; the dated GICS sector history; the link from
Compustat firms (gvkey) to CRSP securities (PERMNO) \\
\bottomrule
\end{tabularx}

\dec{2026-09-18 and 2026-09-26, Q3, brief 06} CRSP and Compustat through the university replace all free data
(yfinance, Stooq, Wikipedia). Wind remains a possible China extension only.

\begin{whybox}
\begin{itemize}
\item \textbf{Free data is survivorship-biased.} Free sources mostly cover firms that exist today; firms that
  were delisted, merged or went bankrupt disappear, and they are exactly the firms with the worst returns.
  A model trained on survivors learns a world that is too optimistic. CRSP keeps every security and its
  delisting return.
\item \textbf{Point-in-time membership.} CRSP records who was in the S\&P 500 on each day; Wikipedia's list is
  today's list.
\item \textbf{Fundamentals with dates.} Compustat records when each quarter's numbers were reported and filed,
  which is what makes look-ahead-free fundamentals possible.
\item \textbf{Comparability.} CRSP and Compustat are the standard data of the literature this thesis builds on
  (Gu, Kelly \& Xiu; Jensen, Kelly \& Pedersen), so results can be set beside theirs.
\end{itemize}
\giveup{nothing material; the cost is licence discipline (no licensed data leave \texttt{Data/}).}
\end{whybox}

\subsection{Sample period}

\dec{2026-10-05, Q16 (a)} \textbf{2000--2024} (previously 2015--2024). Feature history reaches further back (about
1995, for the five-year windows); the \emph{sample} is what is trained and tested on.

\begin{whybox}
\begin{itemize}
\item \textbf{The research question is about regimes, so the binding quantity is the number of stress
  episodes, not the number of days.} 2015--2024 contains essentially one major crisis (COVID). 2000--2024
  contains 2000--02, 2008--09, 2010, 2011, 2015--16, 2018, 2020 and 2022. A study about behaviour at
  regime transitions cannot rest on one transition, and a stress expert cannot specialise if it has seen
  almost no stress.
\item \textbf{2000 is the earliest start with sector data}: Compustat's dated GICS history begins in mid-1999,
  and the industry block and the industry-level gate need it.
\item \textbf{A fairly homogeneous market structure.} Almost all of the sample is after decimalisation (2001),
  so quoted spreads and short-horizon features mean roughly the same thing throughout.
\end{itemize}
\giveup{the pre-2000 episodes (1987, 1998); in exchange, every data source is consistent over the whole
sample.}
\end{whybox}

\subsection{Universe}

\dec{2026-09-26, brief 06} The \textbf{S\&P 500, point in time}, from CRSP membership spells of index 1000500
(index 1000502 has no constituents), bounds inclusive: 502--508 members a day. Dual-class companies are
kept as two securities.

\begin{whybox}
\begin{itemize}
\item \textbf{Short-horizon signals must be tradeable to mean anything.} At a 5-day horizon, effects in
  microcaps are often unimplementable after costs, and their returns are noisy and partly driven by
  bid--ask bounce. Large caps keep the study close to what a fund could actually trade.
\item \textbf{Clean data.} Quotes, open/high/low prices and fundamentals are essentially complete for S\&P 500
  members, so missing-value handling plays a small role.
\item \textbf{Point in time, so no survivorship.} A stock is in the universe exactly on the days it was in the
  index.
\item \textbf{The closest precedent uses large caps too}, which keeps the comparison meaningful.
\end{itemize}
\giveup{breadth. By the fundamental law of active management, performance grows with the square root of the
number of independent bets, so 500 names offer less than 3,000; size and value spreads are also narrower
inside the S\&P 500.}
\end{whybox}

\subsection{Returns and delistings}

\begin{itemize}
\item Daily log return $\log(1+\text{DlyRet})$.
\item \dec{2026-09-26} \textbf{Delisting returns are already in CRSP's daily return} (checked on the data:
  delisting days carry the delisting return), so they are not compounded in again.
\item A forward window that runs past a delisting is completed by a post-delisting fill, default ``cash''
  (zero return). This default is \textbf{not a decision} and is recorded on every run.
\end{itemize}

\begin{whybox}
Leaving out delisting returns biases returns upward, because delistings for poor performance carry large
negative final returns (Shumway 1997). Adding them a second time would bias them downward. Checking the data
showed the CIZ daily return already contains them, so the right rule is to use it as is. Log returns are
used because they add up across days (a 5-day return is the sum of five daily ones) and are closer to
Gaussian, which suits the likelihood.
\end{whybox}

\subsection{Licence and integrity rules}

CRSP and Compustat data, and anything derived per security, stay in \texttt{Data/}, outside version
control; \texttt{results/} holds aggregate outputs only. No out-of-sample performance number on the real
panel (any date from 2010, and the pilot's validation years) is computed by an assistant; pilot and
scored runs are run by Tom, after pre-registration.

\begin{whybox}
The licence forbids redistribution, and a public repository is redistribution. The integrity rule exists
because the comparison is between small effects: one look at test results that feeds back into a choice
contaminates every later number.
\end{whybox}
